\documentclass{article}%
\usepackage[T1]{fontenc}%
\usepackage[utf8]{inputenc}%
\usepackage{lmodern}%
\usepackage{textcomp}%
\usepackage{lastpage}%
\usepackage{geometry}%
\geometry{margin=1in,headheight=14pt,headsep=25pt}%
\usepackage{hyperref}%
\usepackage{enumitem}%
\usepackage{fancyhdr}%
\usepackage{xcolor}%
%

            \hypersetup{
                colorlinks=true,
                linkcolor=blue,
                filecolor=magenta,
                urlcolor=blue,
            }
            
            \pagestyle{fancy}
            \fancyhf{}
            \rhead{Generated on \today}
            \lhead{Course Summary}
            \cfoot{\thepage}
            
            % Configure itemize settings globally
            \setlist[itemize]{nosep}  % Reduce vertical spacing between items
            \setlist[itemize]{leftmargin=*}  % Align with left margin
        %
\title{Linear Programming Course Summary}%
\author{Generated by Scribe.AI}%
\date{\today}%
%
\begin{document}%
\normalsize%
\maketitle%
\section{Key Terms}%
\label{sec:KeyTerms}%
\begin{itemize}[leftmargin=*]%
\subsection{linear programming fundamentals}%
\label{subsec:linearprogrammingfundamentals}%
\begin{itemize}[leftmargin=*]%
%
\end{itemize}

%
\subsection{solution space analysis}%
\label{subsec:solutionspaceanalysis}%
\begin{itemize}[leftmargin=*]%
\item \href{https://www.math.purdue.edu/~yipn/421/2024-08-27-ExSimplex.pdf#page=4}{\textbf{\textcolor{blue}{Feasible Region}}}: \emph{The set of all points that satisfy all the constraints of a linear programming problem.}%
\end{itemize}

%
\subsection{simplex method algorithms}%
\label{subsec:simplexmethodalgorithms}%
\begin{itemize}[leftmargin=*]%
%
\end{itemize}

%
\end{itemize}%
\newpage

%
\section{Problem Types}%
\label{sec:ProblemTypes}%
\begin{itemize}[leftmargin=*]%
\subsection{linear programming fundamentals}%
\label{subsec:linearprogrammingfundamentals}%
\begin{itemize}[leftmargin=*]%
%
\end{itemize}

%
\subsection{advanced linear programming}%
\label{subsec:advancedlinearprogramming}%
\begin{itemize}[leftmargin=*]%
\item \href{https://www.math.purdue.edu/~yipn/421/2024-08-27-ExSimplex.pdf#page=30}{\textbf{\textcolor{blue}{Handling Infeasible Origins}}}: \emph{The origin $(0, 0)$ may not satisfy all constraints in a linear programming problem. Techniques like introducing an auxiliary problem are used to find a feasible starting point.}%
\end{itemize}

%
\end{itemize}%
\newpage

%
\section{Algorithm Solutions}%
\label{sec:AlgorithmSolutions}%
\begin{itemize}[leftmargin=*]%
\subsection{linear optimization techniques}%
\label{subsec:linearoptimizationtechniques}%
\begin{itemize}[leftmargin=*]%
\item \href{https://www.math.purdue.edu/~yipn/421/2024-08-27-ExSimplex.pdf#page=32}{\textbf{\textcolor{blue}{Auxiliary Problem}}}: \emph{When the origin is not feasible, introduce a new variable $x_0$ and modify the constraints to create an auxiliary problem that maximizes $-x_0$. Solve this auxiliary problem using the Simplex method. If the optimal solution has $x_0 = 0$, then the solution is feasible for the original problem.}%
\end{itemize}

%
\end{itemize}%
\newpage

%
\end{document}